\chapter{Intracranial Hemorrhage}

\section*{Patient focus}
This chapter explains the examination or clinical problem in accessible language. Patients should use it to prepare questions and understand the purpose, experience, benefits, and limitations of imaging; local instructions and the official report take priority.


\label{ch:ch76_intracranial_hemorrhage}

Intracranial hemorrhage (ICH) refers to any bleeding within the skull. It can be life-threatening and requires rapid detection. A non-contrast CT (NCCT) of the head is the first-line imaging test in the emergency setting because it is fast, widely available, and highly sensitive for acute blood.

---

\section{1. Patient Perspective: What a Brain CT Shows When There Is Bleeding}

If doctors suspect bleeding in your brain, you will have an \textbf{emergency CT scan} of your head.

\begin{itemize}
  \item You lie on the scanner table, which slides through the machine. The scan takes only 5--10 minutes and does not usually require contrast dye.
  \item On a CT scan, fresh blood appears \textbf{bright white} (hyperdense) against the grey brain tissue. This makes it easy to detect.
  \item The scan shows the \textbf{location} of the bleed (inside the brain tissue, between the brain and skull, or in the fluid spaces around the brain), its \textbf{size}, and whether it is causing pressure on surrounding structures.
  \item If a blood vessel problem is suspected, a second scan with contrast dye (\textbf{CT angiography}) may be performed to look for an aneurysm or arteriovenous malformation.
\end{itemize}

---

\section{2. CT Appearance of Acute Hemorrhage}

Acute blood ($<$ 72 hours) appears hyperdense on NCCT, typically measuring \textbf{50--70 HU}. This is due to the high protein (haemoglobin) content of freshly clotted blood.

\begin{table}[h]
\centering
\begin{tabular}{l l l}
\hline
\textbf{Age of Blood} & \textbf{CT Density} & \textbf{Appearance} \\
\hline
Acute ($<$ 72 h) & 50--70 HU & Bright white (hyperdense) \\
Chronic ($>$ 4 weeks) & CSF density & Fluid-filled cavity \\
Early subacute (3--7 days) & Isodense to brain & May be difficult to detect \\
Late subacute (1--4 weeks) & Hypodense & Darker than brain \\
\hline
\end{tabular}
\caption{CT Appearance of Blood Over Time}
\end{table}

\begin{pearlbox}[The Isodense Trap]
Between approximately 3 and 10 days, subdural hematomas may become isodense with the brain parenchyma, making them very difficult to detect on non-contrast CT. If clinical suspicion remains high despite a normal-appearing NCCT, MRI or contrast-enhanced CT should be considered.
\end{pearlbox}

---

\section{3. Location-Based Differential Diagnosis}

The location of intracranial hemorrhage is the most important factor in determining the underlying cause.

\subsection{Epidural Hematoma}
\begin{itemize}
  \item \textbf{Location}: Between the skull and dura. Typically temporal.
  \item \textbf{Shape}: Biconvex (lens-shaped), does not cross suture lines.
  \item \textbf{Cause}: Usually arterial (middle meningeal artery rupture from temporal bone fracture).
  \item \textbf{Classic presentation}: ``Lucid interval'' --- initial loss of consciousness, then alertness, then rapid deterioration.
  \item \textbf{Management}: Urgent surgical evacuation.
\end{itemize}

\subsection{Subdural Hematoma}
\begin{itemize}
  \item \textbf{Location}: Between dura and arachnoid mater. Usually over cerebral convexity.
  \item \textbf{Shape}: Crescent-shaped, crosses suture lines but not dural attachments (falx, tentorium).
  \item \textbf{Cause}: Bridging vein tears (trauma, especially in elderly, alcoholics, patients on anticoagulants).
  \item \textbf{Chronic subdural}: May present insidiously with confusion, gait disturbance, or focal deficit weeks to months after minor trauma. Appears hypodense on CT.
\end{itemize}

\subsection{Subarachnoid Hemorrhage}
\begin{itemize}
  \item \textbf{Location}: Within the subarachnoid space (CSF-filled spaces around the brain).
  \item \textbf{Appearance}: Hyperdense blood within the sulci, basal cisterns, and fissures.
  \item \textbf{Cause}: Ruptured saccular aneurysm ($\sim$85\%), arteriovenous malformation, or trauma.
  \item \textbf{Classic presentation}: ``Worst headache of my life'' (thunderclap headache).
  \item \textbf{Work-up}: CTA or digital subtraction angiography to identify the source.
\end{itemize}

\subsection{Intraventricular Hemorrhage}
\begin{itemize}
  \item Blood within the lateral, third, or fourth ventricles.
  \item May be primary (e.g., in premature neonates) or secondary extension from intraparenchymal or subarachnoid hemorrhage.
  \item Can cause acute hydrocephalus by obstructing CSF flow.
\end{itemize}

\subsection{Intraparenchymal Hemorrhage}
\begin{itemize}
  \item \textbf{Location}: Within the brain parenchyma itself.
  \item \textbf{Common sites}: Basal ganglia, thalamus, pons, cerebellum (hypertensive hemorrhage); lobar locations (amyloid angiopathy, tumour, vascular malformation).
  \item \textbf{CT}: Hyperdense collection within the parenchyma, often with surrounding hypodense oedema.
  \item \textbf{Volume estimation}: ABC/2 formula (A $\times$ B $\times$ C / 2, where A, B, C are the three perpendicular diameters).
\end{itemize}

---

\section{4. Mass Effect and Herniation Signs}

Intracranial hemorrhage can cause mass effect, displacing brain tissue and potentially leading to fatal herniation:

\begin{itemize}
  \item \textbf{Midline shift}: Measured at the level of the septum pellucidum. Shift $>$ 5 mm is clinically significant and may require surgical intervention.
  \item \textbf{Uncal herniation}: Medial temporal lobe (uncus) herniates through the tentorial notch, compressing the ipsilateral third cranial nerve (fixed, dilated pupil) and the contralateral cerebral peduncle.
  \item \textbf{Subfalcine herniation}: Cingulate gyrus herniates beneath the falx cerebri, potentially compressing the anterior cerebral artery.
  \item \textbf{Tonsillar herniation}: Cerebellar tonsils herniate through the foramen magnum, compressing the brainstem. Often fatal.
  \item \textbf{Effacement of basal cisterns}: Loss of the normal CSF spaces around the brainstem indicates raised intracranial pressure.
\end{itemize}

\begin{pearlbox}[Urgent Findings]
A midline shift $>$ 5 mm, compression of the basal cisterns, or acute hydrocephalus in the setting of intracranial hemorrhage requires immediate communication with the neurosurgical team.
\end{pearlbox}

---

\section{5. When to Recommend CTA: The Spot Sign}

CT angiography (CTA) should be considered in patients with intracranial hemorrhage to identify an underlying vascular lesion:

\begin{itemize}
  \item \textbf{CTA spot sign}: Active contrast extravasation within the hematoma on CTA, indicating ongoing active bleeding. It is a strong predictor of hematoma expansion and is associated with worse outcomes.
  \item \textbf{Indications for CTA}: Lobar hemorrhage in a young patient, subarachnoid hemorrhage, hemorrhage in an unusual location, suspected vascular malformation, or need to assess for aneurysm.
  \item \textbf{Vascular causes to look for}: Saccular aneurysm, arteriovenous malformation, dural arteriovenous fistula, moyamoya, vasculitis, tumour (e.g., glioblastoma, metastasis with haemorrhage).
\end{itemize}

---

\section{6. Structured Reporting Example}

\begin{testbox}[Indications:]
78-year-old male on warfarin for atrial fibrillation. Sudden onset left-sided weakness and decreasing consciousness. Emergency NCCT head requested.
\end{testbox}
\begin{refbox}[Findings:]
\textbf{Hemorrhage}: A hyperdense collection measuring approximately $6 \times 4 \times 3$ cm (estimated volume $\sim$36 mL) is identified in the right basal ganglia, consistent with an acute intraparenchymal hemorrhage. A small amount of blood extends into the right lateral ventricle (intraventricular extension).\\
\textbf{Mass Effect}: There is effacement of the right lateral ventricle and midline shift of 7 mm to the left at the level of the septum pellucidum. The right basal cistern is partially effaced.\\
\textbf{Perilesional Oedema}: Mild hypodensity surrounding the hematoma, consistent with early vasogenic oedema.\\
\textbf{Other}: No acute infarction. No extra-axial hemorrhage. No hydrocephalus. The grey-white matter differentiation is preserved.\\
\textbf{CTA Spot Sign}: Positive --- a focal area of contrast extravasation is seen within the medial aspect of the hematoma, consistent with active bleeding.
\end{refbox}
\begin{pearlbox}[Impression:]
1. Acute right basal ganglia intraparenchymal hemorrhage with intraventricular extension (estimated volume 36 mL).
2. Mass effect with 7 mm midline shift and partial effacement of basal cisterns.
3. Positive CTA spot sign indicating active hemorrhage.
4. Urgent neurosurgical consultation recommended.
\end{pearlbox}

\clearpage
